\section{Processi di Supporto}
\subsection{Documentazione}
Il processo di documentazione si occupa di definire le regole, strumenti e ciclo di vita per la produzione, la revisione, l'approvazione e la pubblicazione della documentazione. Inoltre tutta la documentazione prodotta, sia di tipo interna che di tipo esterna, dovrà seguire il relativo template.
La documentazione deve essere:
\begin{itemize}
	\item Coerente con le Norme di progetto;
	\item Verificabile tramite Pull Request;
	\item Tracciabile rispetto alla modifiche effettuate;
	\item Consultabile nella repository e sito;
\end{itemize}
\subsection{Tipologia Documentazione}
\subsubsection{Documentazione interna}
\begin{itemize}
	\item \textbf{Norme di Progetto};
	\item \textbf{Glossario};
	\item \textbf{Verbali interni};
\end{itemize}

\subsubsection{Documentazione esterna}
\begin{itemize}
	\item \textbf{Verbali Esterni};
	\item \textbf{Diario di Bordo};
\end{itemize}

\subsection{Strumenti per la Documentazione}
La documentazione deve essere prodotta tramite l'ausilio dei seguenti strumenti:
\begin{itemize}
	\item Latex con TexStudios 
\end{itemize}

\subsection{Convenzioni Redazionali}
\begin{itemize}
	\item \textbf{Nome dei file}: tutti i nomi dei file devono essere scritti in minuscolo, utilizzando il carattere \_ come separatore. \\
	Esempio: 
	\item \textbf{Nome del file principale}: il file principale deve essere nominato secondo il nome previsto dal documento finale.
	\item \textbf{Tabelle}: le tabelle di grandi dimensioni devono essere realizzate tramite xltabular o longtable in modo che si possano distribuire automaticamente nello spazio. 
\end{itemize}